\documentclass[10pt,a4paper]{amsart}
\usepackage{amsmath,amssymb,amsthm,graphicx}
% Retained minimal preamble for the Stacks Project chapter more-morphisms.tex (More on Morphisms).
% Verbatim excerpt of the chapter preamble: Xy-pic setup (native preamble lines 14-18) and the native
% theorem declarations with the mathematical macros (native preamble lines 154-217). Copied mechanically;
% nothing added, removed or redefined. The wrapper class and the amsmath/amssymb/amsthm/graphicx base
% packages are supplied by the assembler, so the class/externaldocument part of the native preamble is omitted.
% For commutative diagrams we use Xy-pic
\usepackage[all]{xy}

% We use 2cell for 2-commutative diagrams.
\xyoption{2cell}
% Theorem environments.
%
\theoremstyle{plain}
\newtheorem{theorem}[subsection]{Theorem}
\newtheorem{proposition}[subsection]{Proposition}
\newtheorem{lemma}[subsection]{Lemma}

\theoremstyle{definition}
\newtheorem{definition}[subsection]{Definition}
\newtheorem{example}[subsection]{Example}
\newtheorem{exercise}[subsection]{Exercise}
\newtheorem{situation}[subsection]{Situation}

\theoremstyle{remark}
\newtheorem{remark}[subsection]{Remark}
\newtheorem{remarks}[subsection]{Remarks}

\numberwithin{equation}{subsection}

% Macros
%
\def\lim{\mathop{\mathrm{lim}}\nolimits}
\def\colim{\mathop{\mathrm{colim}}\nolimits}
\def\Spec{\mathop{\mathrm{Spec}}}
\def\Hom{\mathop{\mathrm{Hom}}\nolimits}
\def\Ext{\mathop{\mathrm{Ext}}\nolimits}
\def\SheafHom{\mathop{\mathcal{H}\!\mathit{om}}\nolimits}
\def\SheafExt{\mathop{\mathcal{E}\!\mathit{xt}}\nolimits}
\def\Sch{\mathit{Sch}}
\def\Mor{\mathop{\mathrm{Mor}}\nolimits}
\def\Ob{\mathop{\mathrm{Ob}}\nolimits}
\def\Sh{\mathop{\mathit{Sh}}\nolimits}
\def\NL{\mathop{N\!L}\nolimits}
\def\CH{\mathop{\mathrm{CH}}\nolimits}
\def\proetale{{pro\text{-}\acute{e}tale}}
\def\etale{{\acute{e}tale}}
\def\QCoh{\mathit{QCoh}}
\def\Ker{\mathop{\mathrm{Ker}}}
\def\Im{\mathop{\mathrm{Im}}}
\def\Coker{\mathop{\mathrm{Coker}}}
\def\Coim{\mathop{\mathrm{Coim}}}

% Boxtimes
%
\DeclareMathSymbol{\boxtimes}{\mathbin}{AMSa}{"02}

%
% Macros for moduli stacks/spaces
%
\def\QCohstack{\mathcal{QC}\!\mathit{oh}}
\def\Cohstack{\mathcal{C}\!\mathit{oh}}
\def\Spacesstack{\mathcal{S}\!\mathit{paces}}
\def\Quotfunctor{\mathrm{Quot}}
\def\Hilbfunctor{\mathrm{Hilb}}
\def\Curvesstack{\mathcal{C}\!\mathit{urves}}
\def\Polarizedstack{\mathcal{P}\!\mathit{olarized}}
\def\Complexesstack{\mathcal{C}\!\mathit{omplexes}}
% \Pic is the operator that assigns to X its picard group, usage \Pic(X)
% \Picardstack_{X/B} denotes the Picard stack of X over B
% \Picardfunctor_{X/B} denotes the Picard functor of X over B
\def\Pic{\mathop{\mathrm{Pic}}\nolimits}
\def\Picardstack{\mathcal{P}\!\mathit{ic}}
\def\Picardfunctor{\mathrm{Pic}}
\def\Deformationcategory{\mathcal{D}\!\mathit{ef}}
\usepackage{cleveref}
\crefname{originalreference}{original reference}{original references}
\Crefname{originalreference}{Original reference}{Original references}
\makeatletter
\newlabel{algebra-lemma-strictly-henselian-functorial}{{04GU}{source}{}{sourcepointer-1}{}}
\newlabel{algebra-lemma-strictly-henselian-functorial@cref}{{[originalreference][0][]04GU}{[1][0][]source}{}{}{}}
\makeatother
\crefname{originallemma}{original lemma}{original lemmas}
\Crefname{originallemma}{Original lemma}{Original lemmas}
\makeatletter
\newlabel{lemma-check-weakly-etale-stalks}{{094Q}{source}{}{sourcepointer-2}{}}
\newlabel{lemma-check-weakly-etale-stalks@cref}{{[originallemma][0][]094Q}{[1][0][]source}{}{}{}}
\makeatother
\crefname{originalreference}{original reference}{original references}
\Crefname{originalreference}{Original reference}{Original references}
\makeatletter
\newlabel{more-algebra-lemma-composition-weakly-etale}{{092J}{source}{}{sourcepointer-3}{}}
\newlabel{more-algebra-lemma-composition-weakly-etale@cref}{{[originalreference][0][]092J}{[1][0][]source}{}{}{}}
\makeatother
\crefname{originalreference}{original reference}{original references}
\Crefname{originalreference}{Original reference}{Original references}
\makeatletter
\newlabel{more-algebra-lemma-dumb-properties-henselization}{{07QM}{source}{}{sourcepointer-4}{}}
\newlabel{more-algebra-lemma-dumb-properties-henselization@cref}{{[originalreference][0][]07QM}{[1][0][]source}{}{}{}}
\makeatother
\crefname{originalreference}{original reference}{original references}
\Crefname{originalreference}{Original reference}{Original references}
\makeatletter
\newlabel{more-algebra-lemma-go-down}{{092K}{source}{}{sourcepointer-5}{}}
\newlabel{more-algebra-lemma-go-down@cref}{{[originalreference][0][]092K}{[1][0][]source}{}{}{}}
\makeatother
\crefname{originalreference}{original reference}{original references}
\Crefname{originalreference}{Original reference}{Original references}
\makeatletter
\newlabel{more-algebra-lemma-weakly-etale-permanence}{{092L}{source}{}{sourcepointer-6}{}}
\newlabel{more-algebra-lemma-weakly-etale-permanence@cref}{{[originalreference][0][]092L}{[1][0][]source}{}{}{}}
\makeatother
\crefname{originalreference}{original reference}{original references}
\Crefname{originalreference}{Original reference}{Original references}
\makeatletter
\newlabel{more-algebra-lemma-when-weakly-etale}{{092N}{source}{}{sourcepointer-7}{}}
\newlabel{more-algebra-lemma-when-weakly-etale@cref}{{[originalreference][0][]092N}{[1][0][]source}{}{}{}}
\makeatother
\crefname{originalreference}{original reference}{original references}
\Crefname{originalreference}{Original reference}{Original references}
\makeatletter
\newlabel{more-algebra-theorem-olivier}{{092Z}{source}{}{sourcepointer-8}{}}
\newlabel{more-algebra-theorem-olivier@cref}{{[originalreference][0][]092Z}{[1][0][]source}{}{}{}}
\makeatother
\title{Weak etaleness is equivalent to an isomorphism on strict henselizations}
\author{The Stacks Project Authors}
\date{stacks-project:more-morphisms.tex:local-snapshot-20261011; GNU FDL 1.2 or later}
\begin{document}
\maketitle
\noindent\textit{Editorial note.} Counted span is the author's own proof environment, source lines 19109-19148, immediately following the statement at 19097-19107.\par\medskip
\noindent\textit{Editorial note.} Seven cross-chapter pointers (Algebra, More on Algebra) plus one same-chapter lemma are recorded in assembly.external\_references.\par\medskip
\noindent\textit{Editorial note.} The statement's two conditions are retained verbatim in the statement span.\par\medskip
\noindent\textit{Allowed standard prerequisites (editorial).}\begin{itemize}
\item The stalkwise criterion for weak etaleness. chapter lemma lemma-check-weakly-etale-stalks, cited in the proof. Reduces the global statement to the local rings, which is what the strict henselization test is about.
\item Permanence of weak etaleness under composition and base change. More on Algebra, lemmas more-algebra-lemma-weakly-etale-permanence and more-algebra-lemma-composition-weakly-etale, cited in the proof. Keeps weak etaleness along the ring maps that appear in the henselization diagram.
\item Olivier's theorem on weakly etale maps and henselizations. More on Algebra, theorem more-algebra-theorem-olivier, cited at proof lines 19137-19138. Concludes that the comparison map of strict henselizations is an isomorphism.
\item Functoriality of the strict henselization and go-down. Algebra, lemma algebra-lemma-strictly-henselian-functorial with More on Algebra, lemmas more-algebra-lemma-dumb-properties-henselization and more-algebra-lemma-go-down, cited in the proof. Controls the behaviour of strict henselizations along the maps of local rings.
\end{itemize}\medskip
\ifdefined\hypertarget\hypertarget{sourcepointer-1}{}\fi
\noindent\textit{Source reference.} Original reference 04GU: focused-seeds/geometry/S9-morph-weakly-etale-strictly-henselian/ref-tags.txt, source lines 5--5. Navigation pointer only: the counted proof cites this statement by \textbackslash{}ref as a standing input; the deduction the counted proof has to make is carried out inside the counted span, and no part of that deduction is delegated to this pointer.\par\medskip
\ifdefined\hypertarget\hypertarget{sourcepointer-2}{}\fi
\noindent\textit{Source reference.} Original lemma 094Q: sources/stacks-more-morphisms/pinned-original-source.tex, source lines 18892--18892. Navigation pointer only: the counted proof cites this statement by \textbackslash{}ref as a standing input; the deduction the counted proof has to make is carried out inside the counted span, and no part of that deduction is delegated to this pointer.\par\medskip
\ifdefined\hypertarget\hypertarget{sourcepointer-3}{}\fi
\noindent\textit{Source reference.} Original reference 092J: focused-seeds/geometry/S9-morph-weakly-etale-strictly-henselian/ref-tags.txt, source lines 6--6. Navigation pointer only: the counted proof cites this statement by \textbackslash{}ref as a standing input; the deduction the counted proof has to make is carried out inside the counted span, and no part of that deduction is delegated to this pointer.\par\medskip
\ifdefined\hypertarget\hypertarget{sourcepointer-4}{}\fi
\noindent\textit{Source reference.} Original reference 07QM: focused-seeds/geometry/S9-morph-weakly-etale-strictly-henselian/ref-tags.txt, source lines 7--7. Navigation pointer only: the counted proof cites this statement by \textbackslash{}ref as a standing input; the deduction the counted proof has to make is carried out inside the counted span, and no part of that deduction is delegated to this pointer.\par\medskip
\ifdefined\hypertarget\hypertarget{sourcepointer-5}{}\fi
\noindent\textit{Source reference.} Original reference 092K: focused-seeds/geometry/S9-morph-weakly-etale-strictly-henselian/ref-tags.txt, source lines 8--8. Navigation pointer only: the counted proof cites this statement by \textbackslash{}ref as a standing input; the deduction the counted proof has to make is carried out inside the counted span, and no part of that deduction is delegated to this pointer.\par\medskip
\ifdefined\hypertarget\hypertarget{sourcepointer-6}{}\fi
\noindent\textit{Source reference.} Original reference 092L: focused-seeds/geometry/S9-morph-weakly-etale-strictly-henselian/ref-tags.txt, source lines 9--9. Navigation pointer only: the counted proof cites this statement by \textbackslash{}ref as a standing input; the deduction the counted proof has to make is carried out inside the counted span, and no part of that deduction is delegated to this pointer.\par\medskip
\ifdefined\hypertarget\hypertarget{sourcepointer-7}{}\fi
\noindent\textit{Source reference.} Original reference 092N: focused-seeds/geometry/S9-morph-weakly-etale-strictly-henselian/ref-tags.txt, source lines 10--10. Navigation pointer only: the counted proof cites this statement by \textbackslash{}ref as a standing input; the deduction the counted proof has to make is carried out inside the counted span, and no part of that deduction is delegated to this pointer.\par\medskip
\ifdefined\hypertarget\hypertarget{sourcepointer-8}{}\fi
\noindent\textit{Source reference.} Original reference 092Z: focused-seeds/geometry/S9-morph-weakly-etale-strictly-henselian/ref-tags.txt, source lines 11--11. Navigation pointer only: the counted proof cites this statement by \textbackslash{}ref as a standing input; the deduction the counted proof has to make is carried out inside the counted span, and no part of that deduction is delegated to this pointer.\par\medskip
\begin{lemma}
\label{lemma-weakly-etale-strictly-henselian-local-rings}
Let $f : X \to Y$ be a morphism of schemes.
The following are equivalent
\begin{enumerate}
\item $f$ is weakly \'etale, and
\item for $x \in X$ the local ring map
$\mathcal{O}_{Y, f(x)} \to \mathcal{O}_{X, x}$ induces an isomorphism
on strict henselizations.
\end{enumerate}
\end{lemma}

\begin{proof}
Let $x \in X$ be a point with image $y = f(x)$ in $Y$.
Choose a separable algebraic closure $\kappa^{sep}$ of $\kappa(x)$.
Let $\mathcal{O}_{X, x}^{sh}$ be the strict henselization
corresponding to $\kappa^{sep}$ and $\mathcal{O}_{Y, y}^{sh}$
the strict henselization relative to the separable algebraic
closure of $\kappa(y)$ in $\kappa^{sep}$.
Consider the commutative diagram
$$
\xymatrix{
\mathcal{O}_{X, x} \ar[r] & \mathcal{O}_{X, x}^{sh} \\
\mathcal{O}_{Y, y} \ar[u] \ar[r] & \mathcal{O}_{Y, y}^{sh} \ar[u]
}
$$
local homomorphisms of local rings, see
Algebra, Lemma \ref{algebra-lemma-strictly-henselian-functorial}.
Since the strict henselization is a filtered colimit of \'etale
ring maps, More on Algebra, Lemma \ref{more-algebra-lemma-when-weakly-etale}
shows the horizontal maps are weakly \'etale.
Moreover, the horizontal maps are faithfully flat by
More on Algebra, Lemma \ref{more-algebra-lemma-dumb-properties-henselization}.

\medskip\noindent
Assume $f$ weakly \'etale. By Lemma \ref{lemma-check-weakly-etale-stalks}
the left vertical arrow is weakly \'etale. By
More on Algebra, Lemmas \ref{more-algebra-lemma-composition-weakly-etale} and
\ref{more-algebra-lemma-weakly-etale-permanence}
the right vertical arrow is weakly \'etale. By
More on Algebra, Theorem \ref{more-algebra-theorem-olivier}
we conclude the right vertical map is an isomorphism.

\medskip\noindent
Assume $\mathcal{O}_{Y, y}^{sh} \to \mathcal{O}_{X, x}^{sh}$ is an isomorphism.
Then $\mathcal{O}_{Y, y} \to \mathcal{O}_{X, x}^{sh}$ is weakly \'etale.
Since $\mathcal{O}_{X, x} \to \mathcal{O}_{X, x}^{sh}$ is faithfully
flat we conclude that $\mathcal{O}_{Y, y} \to \mathcal{O}_{X, x}$
is weakly \'etale by
More on Algebra, Lemma \ref{more-algebra-lemma-go-down}.
Thus (2) implies (1) by Lemma \ref{lemma-check-weakly-etale-stalks}.
\end{proof}

\end{document}
